%\documentclass[DiscStateSpaceToolkit.tex]{subfiles} 
%\begin{document}

\cite{Imrohoroglu1989} introduces the use of both idiosyncratic and aggregate shocks; albeit without an aggregate shock.
The paper looks at the cost of recessions (business cycles) in a manner following \cite{Lucas1987} introducing into the model
the idea that recessions do not affect everyone in the same way, in particular in this model in a recession those how
find themselves out of work will be more likely to remain out of work during a recession. A number of different economies
(model variants) are considered and the code has an option at the beginning for which one will be implemented.
It is of additional interest here as it will illustrate how to use the toolkit commands in the presence of two exogenous shocks.

The value function problem to be solved is,
\begin{equation*}
 V(a,s,z)=\sup_{a'} \{ \frac{c^{1-\sigma}}{1-\sigma}+\beta*E[V(a',s',z')|s,z]
\end{equation*}
subject to
\begin{align*}
 & c+\frac{a'}{1+r}=a+earnings(s) \\
 & earnings(s)=\left\{ \begin{array}{lr} y & \text{if s=employed}, \\ \theta y & \text{if s=unemployed} \end{array} \right.
 & (s,z')=\pi(s,z)
\end{align*}
where $a$ is asset holdings, $c$ is consumption, $s$ is an idiosyncratic shock and determines whether
a household is employed or unemployed. $z$ is an aggregate shock (recession or expansion). 
A variety of different economies are considered, with and without aggregate shocks, $z$, and changing
the way in which $s$ changes state (whether or not it depends on aggregate state). See \cite{Imrohoroglu1989}
for a detailed discussion.


The following code solves this model, using parallelization on GPU, and (IN PROGRESS) replicates the results
of \cite{Imrohoroglu1989}; it contains comments describing what is being done
at each step. This code is a copy of $Imrohoroglu1989.m$ which you can run and modify to get a better feel for how
to use the toolkit. Below that is the code that defines the return function, $Imrohoroglu1989\_{}ReturnFn.m$.

\textbf{Contents of Imrohoroglu1989.m}
\lstinputlisting[language=Matlab]{./Examples/MatlabCodes/Imrohoroglu1989.m}

\smallskip

\textbf{Contents of Imrohoroglu1989$\_{}$ReturnFn.m}
\lstinputlisting[language=Matlab]{./Examples/MatlabCodes/Imrohoroglu1989_ReturnFn.m}

%\end{document}